\documentclass[12pt,a4paper]{article}

% --- Encodage & langue ---
\usepackage[utf8]{inputenc}
\usepackage[T1]{fontenc}
\usepackage[french]{babel}

% --- Mise en page ---
\usepackage[top=2.5cm, bottom=2.5cm, left=2.5cm, right=2.5cm]{geometry}
\usepackage{setspace}
\onehalfspacing
\usepackage{parskip}
\setlength{\parindent}{0pt}

% --- Typographie ---
\usepackage{lmodern}
\usepackage{microtype}

% --- Couleurs & graphiques ---
\usepackage[dvipsnames]{xcolor}
\usepackage{graphicx}
\usepackage{tikz}
\usetikzlibrary{shapes.geometric, arrows.meta, positioning, fit, backgrounds, calc}
\usepackage{tcolorbox}
\tcbuselibrary{skins, breakable}

% --- Tableaux ---
\usepackage{booktabs}
\usepackage{tabularx}
\usepackage{array}
\usepackage{multirow}
\usepackage{longtable}

% --- Code ---
\usepackage{listings}
\usepackage{fancyvrb}

\definecolor{codebg}{RGB}{248,248,248}
\definecolor{codeframe}{RGB}{200,200,200}
\definecolor{keywordcolor}{RGB}{0,0,180}
\definecolor{commentcolor}{RGB}{100,150,100}
\definecolor{stringcolor}{RGB}{180,0,0}

\lstset{
  language=Python,
  backgroundcolor=\color{codebg},
  frame=single,
  rulecolor=\color{codeframe},
  basicstyle=\ttfamily\footnotesize,
  keywordstyle=\color{keywordcolor}\bfseries,
  commentstyle=\color{commentcolor}\itshape,
  stringstyle=\color{stringcolor},
  showstringspaces=false,
  breaklines=true,
  numbers=left,
  numberstyle=\tiny\color{gray},
  numbersep=5pt,
  tabsize=4,
  captionpos=b,
}

% --- Liens & PDF ---
\usepackage[hidelinks]{hyperref}
\usepackage{nameref}

% --- En-têtes & pieds de page ---
\usepackage{fancyhdr}
\pagestyle{fancy}
\fancyhf{}
\renewcommand{\headrulewidth}{0.4pt}
\fancyhead[L]{\small Projet de 1\iere{} année --- Groupe 53}
\fancyhead[R]{\small ENSAI 2025--2026}
\fancyfoot[C]{\thepage}

% --- Titres de sections ---
\usepackage{titlesec}
\titleformat{\section}{\large\bfseries}{}{0em}{\thesection\quad}
\titleformat{\subsection}{\normalsize\bfseries}{}{0em}{\thesubsection\quad}
\titleformat{\subsubsection}{\normalsize\itshape\bfseries}{}{0em}{\thesubsubsection\quad}

% --- Macros utilitaires ---
\newcommand{\sport}[1]{\textbf{#1}}
\newcommand{\code}[1]{\texttt{#1}}
\newcommand{\fichier}[1]{\texttt{#1}}
\newcommand{\classe}[1]{\texttt{\textbf{#1}}}

% Boîte de capture d'écran fictive
\newcommand{\screenshot}[2]{%
  \begin{center}
    \fbox{%
      \begin{minipage}{0.88\textwidth}
        \centering\vspace{0.4cm}
        \textit{[Capture d'écran : #1]}\\[0.3cm]
        \textcolor{gray}{\small #2}
        \vspace{0.4cm}
      \end{minipage}%
    }
  \end{center}
}

% Boîte d'encadré coloré
\newtcolorbox{infobox}[1][]{
  colback=blue!5!white,
  colframe=blue!40!white,
  fonttitle=\bfseries,
  title={#1},
  breakable
}

\begin{document}

% ============================================================
% PAGE DE GARDE
% ============================================================
\begin{titlepage}
  \centering
  \vspace*{1.5cm}

  {\LARGE\bfseries École Nationale de la Statistique\\
  et de l'Analyse de l'Information}\\[0.3cm]
  {\large ENSAI --- Bruz}\\[1.5cm]

  \rule{\linewidth}{1.5pt}\\[0.4cm]
  {\Huge\bfseries Rapport de Projet}\\[0.3cm]
  {\LARGE Traitement de données sportives}\\[0.2cm]
  \rule{\linewidth}{1.5pt}\\[1.2cm]

  {\large\itshape Gestion multi-sports des résultats de compétitions\\
  Application en ligne de commande --- Python}\\[2cm]

  \begin{tabular}{ll}
    \textbf{Groupe :} & 53 \\[0.3cm]
    \textbf{Étudiants :} & Jade \textsc{Colin} \\
                         & Mathias \textsc{Douillard} \\
                         & Adam \textsc{[Nom]} \\
                         & Mohamed Fouad \textsc{Chaabane} \\[0.3cm]
    \textbf{Encadrant :} & Clément \textsc{Valot} \\[0.3cm]
    \textbf{Année académique :} & 2025--2026 \\
  \end{tabular}

  \vfill

  {\small Projet de traitement de données --- 1\iere{} année\\
  Mai 2026}
\end{titlepage}

% ============================================================
% RÉSUMÉ
% ============================================================
\thispagestyle{empty}
\section*{Résumé}

Dans le contexte des Jeux Olympiques de Paris 2024 et de la dynamique sportive
nationale qui en résulte, le Ministère des Sports (client fictif) a exprimé le
besoin d'une application unifiée pour gérer et analyser les résultats de
compétitions sportives. Ce projet de première année à l'ENSAI répond à cette
demande en développant une application en ligne de commande entièrement écrite
en Python.

L'application couvre \textbf{neuf disciplines sportives} : basketball (NBA),
football (Ligue des Champions UEFA et ligues européennes), tennis (circuits ATP
et WTA), League of Legends (EMEA LEC), Counter-Strike~2 (CS2), volleyball
(hommes et femmes), StarCraft~II, badminton et échecs. Pour chaque sport, elle
propose un ensemble de fonctionnalités statistiques — classements, fiches
joueurs, rosters, agendas — accessibles via une interface interactive
paginée.

Du point de vue technique, le projet est structuré en quatre couches
logicielles indépendantes : les classes métier (\code{Model}), les chargeurs
de données (\code{Parsers}), les modules d'analyse (\code{Analysis}) et les
utilitaires communs (\code{Common}). Une couverture de tests unitaires assure
la fiabilité des composants clés. L'ensemble du code respecte les conventions
PEP~8 (formatage \code{black}, linter \code{Flake8}) et est documenté avec des
\textit{docstrings} au format NumPy.

\bigskip
\textbf{Mots-clés :} Python, pandas, traitement de données, sports, CLI, programmation
orientée objet, tests unitaires.

\newpage

% ============================================================
% TABLE DES MATIÈRES
% ============================================================
\thispagestyle{empty}
\tableofcontents
\newpage

% ============================================================
% CORPS DU RAPPORT
% ============================================================
\setcounter{page}{1}

% ------------------------------------------------------------
\section{Introduction}
% ------------------------------------------------------------

\subsection{Contexte général}

Les Jeux Olympiques de Paris 2024 ont renforcé l'intérêt de la sphère publique
française pour la gestion et l'analyse de données sportives. Dans ce contexte,
le \textbf{Ministère des Sports} (client fictif de ce projet pédagogique) a
manifesté le souhait de disposer d'un outil informatique permettant de
centraliser les résultats de compétitions issues de sources hétérogènes, d'en
extraire des statistiques pertinentes et de les restituer de manière claire.

Ce projet de première année à l'ENSAI s'inscrit dans ce cadre en proposant une
application de traitement de données sportives développée intégralement en
Python. Le groupe 53 --- composé de quatre étudiants encadrés par
Clément~\textsc{Valot} --- a mené ce projet sur l'année académique 2025--2026.

\subsection{Objectifs du projet}

L'application doit permettre :
\begin{itemize}
  \item d'ingérer et de nettoyer des jeux de données hétérogènes au format CSV
        couvrant plusieurs sports ;
  \item de proposer des analyses statistiques adaptées à chaque discipline ;
  \item d'offrir une interface utilisateur interactive en ligne de commande ;
  \item de respecter les bonnes pratiques du génie logiciel (tests, documentation,
        versionnage, style de code).
\end{itemize}

\subsection{Structure du rapport}

Le présent rapport décrit d'abord le cahier des charges et la répartition du
travail au sein du groupe (sections~\ref{sec:cdc} et~\ref{sec:roles}), puis
présente les données utilisées (section~\ref{sec:donnees}), l'architecture
logicielle (sections~\ref{sec:archi} et~\ref{sec:model}), les fonctionnalités
implémentées (section~\ref{sec:fonct}), la démarche de tests
(section~\ref{sec:tests}), et enfin quelques résultats illustrés par des
captures d'écran (section~\ref{sec:resultats}).

% ------------------------------------------------------------
\section{Cahier des charges}
\label{sec:cdc}
% ------------------------------------------------------------

\subsection{Exigences fonctionnelles}

Le cahier des charges impose de couvrir \textbf{au minimum cinq sports
différents} avec, pour chacun, au moins trois fonctionnalités distinctes
(classements, fiches joueurs, statistiques d'équipes, etc.). L'application
doit également intégrer un \textbf{agenda commun} listant les derniers matchs
toutes disciplines confondues.

Les fonctionnalités retenues pour l'ensemble de l'application sont les
suivantes :

\begin{center}
\begin{tabularx}{\textwidth}{lX}
  \toprule
  \textbf{Catégorie} & \textbf{Fonctionnalité} \\
  \midrule
  Général & Agenda des derniers matchs (toutes disciplines) \\
  Général & Export CSV des résultats affichés \\
  Équipe  & Statistiques moyennes par match \\
  Équipe  & Roster (liste des joueurs et coachs) \\
  Équipe  & Classement par points \\
  Joueurs & Fiche individuelle complète \\
  Joueurs & Liste paginée avec filtres \\
  Compétition & Bracket de tournoi \\
  Sport spécifique & Cf.\ tableau section~\ref{sec:fonct} \\
  \bottomrule
\end{tabularx}
\end{center}

\subsection{Exigences techniques}

Le projet est soumis aux contraintes techniques suivantes, imposées par les
enseignants :

\begin{itemize}
  \item \textbf{Langage} : Python uniquement (aucune interface graphique).
  \item \textbf{Bibliothèques autorisées} : pandas (manipulation de données),
        pytest et pytest-cov (tests), black (formatage), Flake8 (linter).
  \item \textbf{Versionnage} : dépôt GitHub privé, avec l'encadrant ajouté
        comme collaborateur.
  \item \textbf{Documentation} : \textit{docstrings} au format NumPy ou Google
        pour toutes les fonctions et classes publiques.
  \item \textbf{Tests} : couverture par tests unitaires avec pytest ; rapport
        de couverture via pytest-cov.
  \item \textbf{Structure} : fichier \fichier{README.md}, fichier
        \fichier{requirements.txt} et point d'entrée \fichier{\_\_main\_\_.py}.
\end{itemize}

\subsection{Diagramme des cas d'utilisation}

La figure~\ref{fig:usecase} présente les principaux cas d'utilisation de
l'application tels qu'ils ont été identifiés lors de la phase de conception.

\begin{figure}[h!]
\centering
\begin{tikzpicture}[
  actor/.style={draw, circle, minimum size=1cm, font=\small},
  usecase/.style={draw, ellipse, minimum width=3.5cm, minimum height=0.8cm,
                  font=\small, align=center},
  sys/.style={draw, rectangle, dashed, minimum width=9cm, minimum height=10cm,
              font=\bfseries\small},
  >=Stealth
]
  % Acteur
  \node[actor] (user) at (-5.5, 0) {};
  \node[below=0.1cm of user, font=\small] {Utilisateur};
  \draw (-5.5, 0.5) -- (-5.5, 1) -- (-5.2, 0.6);
  \draw (-5.5, 1) -- (-5.8, 0.6);
  \draw (-5.5, 0) -- (-5.5, -0.7);
  \draw (-5.5, -0.7) -- (-5.2, -1.2);
  \draw (-5.5, -0.7) -- (-5.8, -1.2);

  % Système
  \node[sys, anchor=center] (sys) at (1.5, 0) {};
  \node[above=4.7cm of sys, xshift=0cm, font=\bfseries\small] {Application sportive};

  % Cas d'utilisation
  \node[usecase] (agenda)  at (1.5,  3.5) {Consulter l'agenda\\des matchs};
  \node[usecase] (stats)   at (1.5,  2.0) {Voir les stats\\d'une équipe};
  \node[usecase] (roster)  at (1.5,  0.7) {Consulter un roster};
  \node[usecase] (joueur)  at (1.5, -0.7) {Fiche d'un joueur};
  \node[usecase] (classmt) at (1.5, -2.0) {Afficher un classement};
  \node[usecase] (bracket) at (1.5, -3.3) {Voir un bracket};
  \node[usecase] (export)  at (1.5, -4.6) {Exporter en CSV};

  % Flèches
  \foreach \uc in {agenda, stats, roster, joueur, classmt, bracket, export}
    \draw[->, thick] (user) -- (\uc);
\end{tikzpicture}
\caption{Diagramme des cas d'utilisation de l'application.}
\label{fig:usecase}
\end{figure}

% ------------------------------------------------------------
\section{Répartition des rôles}
\label{sec:roles}
% ------------------------------------------------------------

Le groupe a organisé son travail en s'appuyant sur les compétences et les
centres d'intérêt de chacun. Le tableau ci-dessous synthétise la répartition
retenue.

\begin{center}
\begin{tabularx}{\textwidth}{lX}
  \toprule
  \textbf{Membre} & \textbf{Contributions principales} \\
  \midrule
  Jade \textsc{Colin}
    & Conception et implémentation des classes du modèle objet
      (\classe{Joueur}, \classe{Joueurs}, \classe{Equipe}, \classe{Equipes}) ;
      modules d'analyse Chess et Volleyball ; rédaction de tests unitaires
      pour les classes du modèle. \\[0.3cm]
  Adam \textsc{[Nom]}
    & Classe \classe{Match}/\classe{Matchs} et modélisation des compétitions ;
      modules d'analyse Basketball et Football (Ligue des Champions) ;
      contribution aux tests d'intégration. \\[0.3cm]
  Mohamed Fouad \textsc{Chaabane}
    & Interface utilisateur CLI (\fichier{\_\_main\_\_.py}) ;
      modules d'analyse Football (ligues européennes), Tennis ATP/WTA ;
      conception de la navigation paginée et du filtre de recherche. \\[0.3cm]
  Mathias \textsc{Douillard}
    & Architecture générale ; classe \classe{BaseLoader} et tous les loaders
      (\fichier{src/Parsers/}) ; modules d'analyse LoL, CS2, StarCraft~II,
      Badminton ; module \fichier{générique.py} (fonctions partagées) ;
      intégration finale et documentation. \\
  \bottomrule
\end{tabularx}
\end{center}

La gestion du projet a été assurée conjointement via un dépôt GitHub privé,
avec des \textit{pull requests} pour les contributions majeures et des réunions
hebdomadaires de suivi avec l'encadrant Clément~\textsc{Valot}.

% ------------------------------------------------------------
\section{Données et jeux de données}
\label{sec:donnees}
% ------------------------------------------------------------

\subsection{Sources et volumétrie}

L'ensemble des données est stocké dans le répertoire \fichier{data/} et fourni
au format CSV. Chaque sport possède son propre sous-dossier. Le
tableau~\ref{tab:datasets} présente les jeux de données intégrés.

\begin{center}
\begin{longtable}{llll}
  \caption{Jeux de données utilisés dans l'application.}
  \label{tab:datasets}\\
  \toprule
  \textbf{Sport} & \textbf{Fichiers principaux} & \textbf{Saison / Source} \\
  \midrule
  \endfirsthead
  \toprule
  \textbf{Sport} & \textbf{Fichiers principaux} & \textbf{Saison / Source} \\
  \midrule
  \endhead
  \bottomrule
  \endfoot
  Basketball (NBA) & \code{game.csv}, \code{player.csv}, \code{team.csv}
    & Saison 2022--2023 \\
  Football UCL & \code{match.csv}, \code{player.csv}, \code{team.csv}
    & Ligue des Champions UEFA \\
  Football ligues & \code{match.csv}, \code{league.csv}, \code{team.csv}
    & Ligues européennes \\
  League of Legends & \code{match.csv}, \code{player.csv}, \code{coach.csv}, \code{team.csv}
    & EMEA LEC \\
  CS2 & \code{match.csv}, \code{player.csv}, \code{coach.csv}, \code{team.csv}
    & Compétitions professionnelles \\
  Tennis & \code{atp\_matches\_2024.csv}, \code{wta\_matches\_2024.csv}
    & Saison 2024 ATP/WTA \\
  Volleyball & \code{match\_men.csv}, \code{match\_women.csv}, \code{player\_men.csv}
    & Compétitions internationales \\
  StarCraft~II & \code{match.csv}, \code{player.csv}
    & Compétitions professionnelles \\
  Badminton & \code{match.csv}, \code{player.csv}
    & Compétitions internationales \\
  Échecs & \code{match.csv}, \code{player.csv}
    & Tournois classiques \\
\end{longtable}
\end{center}

\subsection{Nettoyage et transformation des données}

Les données brutes nécessitent plusieurs transformations avant d'être
exploitables. Ces opérations sont centralisées dans les loaders
(\fichier{src/Parsers/}).

\subsubsection{Taille des joueurs NBA}

Les tailles dans le jeu de données Basketball sont au format américain
\textit{pieds-pouces} (par exemple \code{"6-8"}). Elles sont converties en
centimètres selon la formule :

\begin{lstlisting}[caption={Conversion taille NBA (BasketballLoader.py)}]
def _height_to_cm(h: str) -> float | None:
    """Convertit '6-8' (pieds-pouces) en centimetres."""
    try:
        feet, inches = map(int, h.split("-"))
        return round((feet * 12 + inches) * 2.54, 1)
    except (ValueError, AttributeError):
        return None
\end{lstlisting}

\subsubsection{Parsing des dates}

Toutes les dates sont standardisées en \code{pd.Timestamp} via la méthode
\code{\_load\_csv()} de \classe{BaseLoader}, qui accepte un paramètre optionnel
\code{date\_cols} :

\begin{lstlisting}[caption={Parsing des dates dans BaseLoader.py}]
for col in (date_cols or []):
    if col in df.columns:
        df[col] = pd.to_datetime(df[col], errors="coerce")
\end{lstlisting}

\subsubsection{Gestion des valeurs manquantes}

Pour les colonnes numériques des statistiques (kills, gold, assists, etc.),
les valeurs manquantes sont implicitement gérées par les méthodes \code{.mean()}
et \code{.fillna(0)} de pandas, selon le contexte. Pour les données textuelles
(noms de joueurs, équipes), la valeur \code{NaN} est préservée et filtrée lors
de l'affichage.

\subsection{Construction du nom complet (Basketball)}

Le jeu de données NBA contient séparément le prénom (\code{first\_name}) et le
nom (\code{last\_name}). Un champ synthétique \code{full\_name} est créé lors du
chargement :

\begin{lstlisting}[caption={Création du full\_name (BasketballLoader.py)}]
players["full_name"] = (
    players["first_name"].str.strip()
    + " "
    + players["last_name"].str.strip()
)
\end{lstlisting}

% ------------------------------------------------------------
\section{Architecture du code}
\label{sec:archi}
% ------------------------------------------------------------

\subsection{Vue d'ensemble}

L'application est structurée selon une architecture en \textbf{quatre couches}
indépendantes, auxquelles s'ajoute le point d'entrée \fichier{\_\_main\_\_.py}.
Cette séparation des responsabilités facilite la maintenance et l'extension
à de nouveaux sports.

\begin{figure}[h!]
\centering
\begin{tikzpicture}[
  layer/.style={draw, fill=#1, rounded corners=4pt,
                minimum width=10cm, minimum height=1.1cm,
                font=\bfseries\small, align=center},
  arr/.style={->, thick, >=Stealth},
  note/.style={font=\footnotesize\itshape, text width=4cm, align=left}
]
  \node[layer=orange!20] (main)    at (0,  5.5) {\_\_main\_\_.py \quad (Interface CLI)};
  \node[layer=blue!15]   (analysis)at (0,  3.8) {src/Analysis/ \quad (Fonctions d'analyse)};
  \node[layer=green!15]  (parsers) at (0,  2.1) {src/Parsers/ \quad (Chargeurs CSV)};
  \node[layer=red!15]    (model)   at (0,  0.4) {src/Model/ \quad (Classes métier)};
  \node[layer=gray!15]   (common)  at (0, -1.2) {src/Common/ \quad (Utilitaires partagés)};

  \draw[arr] (main)     -- (analysis);
  \draw[arr] (analysis) -- (parsers);
  \draw[arr] (parsers)  -- (model);
  \draw[arr] (analysis.west) to[out=180, in=180] (common.west);
  \draw[arr] (parsers.east)  to[out=0,   in=0]   (model.east);

  % Étiquettes
  \node[note] at (7, 4.65)  {stats, classements, rosters, brackets};
  \node[note] at (7, 2.95)  {lecture CSV, nettoyage, jointures};
  \node[note] at (7, 1.25)  {Joueur, Match, Equipe, Sport};
  \node[note] at (7, -0.35) {décorateurs, helpers};
\end{tikzpicture}
\caption{Architecture en couches de l'application.}
\label{fig:archi}
\end{figure}

\subsection{Description des couches}

\subsubsection{src/Model/ --- Classes métier}

Cette couche définit les abstractions orientées objet génériques, indépendantes
de tout sport particulier :

\begin{itemize}
  \item \classe{Joueur} / \classe{Joueurs} : représentation d'un joueur
        individuel et d'une collection de joueurs.
  \item \classe{Match} / \classe{Matchs} : représentation d'un match et
        d'une collection de matchs avec calcul automatique du gagnant.
  \item \classe{Equipe} / \classe{Equipes} : représentation d'une équipe
        et d'une collection d'équipes.
  \item \classe{Sport}, \classe{StatsSport} : méta-informations sur un sport.
\end{itemize}

Les classes sont \textbf{configurables par injection de noms de colonnes} à
l'initialisation, ce qui les rend réutilisables pour tous les sports sans
modification.

\subsubsection{src/Parsers/ --- Chargeurs de données}

Cette couche est basée sur une classe abstraite \classe{BaseLoader} qui
centralise la logique de lecture CSV. Chaque sport possède un loader dédié
héritant de \classe{BaseLoader} et déclarant uniquement son
\code{SPORT\_FOLDER} :

\begin{lstlisting}[caption={Structure minimale d'un loader (exemple Badminton)}]
class BadmintonLoader(BaseLoader):
    SPORT_FOLDER = "badminton"

    def load_all(self):
        players = self._load_csv("player.csv")
        matches = self._load_csv("match.csv", date_cols=["date"])
        return players, matches
\end{lstlisting}

Les loaders disponibles sont : \code{BasketballLoader}, \code{FootballLoader},
\code{FootballChampionsLeagueLoader}, \code{LolLoader}, \code{Cs2Loader},
\code{TennisLoader}, \code{VolleyballLoader}, \code{Starcraft2Loader},
\code{BadmintonLoader}, \code{ChessLoader}.

\subsubsection{src/Analysis/ --- Modules d'analyse}

Chaque sport dispose d'un module Python dédié (par exemple \fichier{basketball.py},
\fichier{tennis.py}) qui expose des fonctions métier de haut niveau. Ces
modules utilisent un \textbf{pattern singleton de chargement paresseux}
(\textit{lazy loading}) : les données ne sont chargées qu'au premier appel et
mises en cache pour les appels suivants :

\begin{lstlisting}[caption={Pattern de chargement paresseux (extrait de basketball.py)}]
_loader = None
_players: pd.DataFrame = None
_teams: pd.DataFrame = None
_matches: pd.DataFrame = None

def _load():
    global _loader, _players, _teams, _matches
    if _loader is None:
        _loader = BasketballLoader()
        _players, _teams, _matches = _loader.load_all()
\end{lstlisting}

Le module \fichier{générique.py} regroupe les fonctions partagées par tous les
sports : \code{formater\_roster}, \code{afficher\_classement},
\code{afficher\_bracket}, \code{fiche\_joueur}, \code{lister\_joueurs},
\code{agenda\_recents}.

\subsubsection{src/Common/ --- Utilitaires}

Ce module regroupe des outils transverses :
\begin{itemize}
  \item \code{print\_timings} : décorateur mesurant le temps d'exécution
        CPU d'une fonction.
  \item \code{parse\_boolean} : conversion robuste d'une chaîne en booléen
        (accepte \code{"true"}, \code{"vrai"}, \code{True}, etc.).
\end{itemize}

\subsubsection{\_\_main\_\_.py --- Interface CLI}

Point d'entrée de l'application, ce fichier (~1\,400 lignes) implémente
l'intégralité de l'interface en ligne de commande. Il s'appuie sur une
configuration déclarative par sport (\code{\_make\_sports\_config()}) qui
associe à chaque sport ses fonctionnalités et leurs paramètres d'entrée.

\subsection{Diagramme de packages}

\begin{figure}[h!]
\centering
\begin{tikzpicture}[
  pkg/.style={draw, fill=yellow!10, font=\small\ttfamily,
              minimum width=2.8cm, minimum height=0.8cm, align=center},
  arr/.style={->, thick, >=Stealth, dashed}
]
  \node[pkg] (main)     at (0,    4)   {\_\_main\_\_.py};
  \node[pkg] (analysis) at (-3,   2)   {src.Analysis};
  \node[pkg] (parsers)  at (0,    2)   {src.Parsers};
  \node[pkg] (model)    at (3,    2)   {src.Model};
  \node[pkg] (common)   at (-3,   0)   {src.Common};
  \node[pkg] (data)     at (0,    0)   {data/};

  \draw[arr] (main) -- (analysis);
  \draw[arr] (main) -- (parsers);
  \draw[arr] (analysis) -- (parsers);
  \draw[arr] (analysis) -- (common);
  \draw[arr] (parsers) -- (model);
  \draw[arr] (parsers) -- (data);
\end{tikzpicture}
\caption{Diagramme de packages de l'application.}
\label{fig:packages}
\end{figure}

% ------------------------------------------------------------
\section{Modélisation orientée objet}
\label{sec:model}
% ------------------------------------------------------------

\subsection{Principes de conception}

Les classes du modèle suivent deux principes fondamentaux :

\begin{enumerate}
  \item \textbf{Généricité par configuration} : plutôt que de sous-classer pour
        chaque sport, les classes acceptent en paramètre les noms des colonnes
        correspondant à leurs attributs (\code{col\_nom}, \code{col\_id}, etc.).
        Cela évite l'explosion du nombre de classes.
  \item \textbf{Encapsulation des DataFrames} : les classes \textit{wrappent}
        des \code{pd.DataFrame} et \code{pd.Series} plutôt que de les remplacer,
        préservant la puissance de pandas tout en exposant une API orientée
        objet cohérente.
\end{enumerate}

\subsection{Diagramme de classes}

\begin{figure}[h!]
\centering
\begin{tikzpicture}[
  cls/.style={draw, fill=white, font=\footnotesize\ttfamily,
              minimum width=3.2cm, align=center},
  inh/.style={->, >=open triangle 90, thick},
  agg/.style={->, >=diamond, thick},
  >=Stealth
]
  % Joueur
  \node[cls, minimum height=2.5cm] (joueur) at (-4, 0) {
    \textbf{Joueur}\\[2pt]
    \rule{3cm}{0.3pt}\\
    id, nom, naissance\\age, equipe, taille\\[2pt]
    \rule{3cm}{0.3pt}\\
    \_\_getitem\_\_()\\to\_dict()
  };

  % Joueurs
  \node[cls, minimum height=2.5cm] (joueurs) at (-4, -4.5) {
    \textbf{Joueurs}\\[2pt]
    \rule{3cm}{0.3pt}\\
    \_df : DataFrame\\[2pt]
    \rule{3cm}{0.3pt}\\
    get\_par\_id/nom/equipe()\\filtrer(), ages()\\age\_moyen(), resume()
  };

  % Match
  \node[cls, minimum height=2.5cm] (match) at (0.5, 0) {
    \textbf{Match}\\[2pt]
    \rule{3cm}{0.3pt}\\
    id, date, equipe\_1/2\\score\_1/2, gagnant\\[2pt]
    \rule{3cm}{0.3pt}\\
    \_\_getitem\_\_()\\to\_dict()
  };

  % Matchs
  \node[cls, minimum height=2.5cm] (matchs) at (0.5, -4.5) {
    \textbf{Matchs}\\[2pt]
    \rule{3cm}{0.3pt}\\
    \_df : DataFrame\\[2pt]
    \rule{3cm}{0.3pt}\\
    get\_par\_equipe/date()\\victoires(), bilan()\\classement(), score\_moyen()
  };

  % Equipe
  \node[cls, minimum height=2.5cm] (equipe) at (5, 0) {
    \textbf{Equipe}\\[2pt]
    \rule{3cm}{0.3pt}\\
    id, nom\\abreviation, ville\\[2pt]
    \rule{3cm}{0.3pt}\\
    get\_joueurs(joueurs)
  };

  % Equipes
  \node[cls, minimum height=2.5cm] (equipes) at (5, -4.5) {
    \textbf{Equipes}\\[2pt]
    \rule{3cm}{0.3pt}\\
    \_df : DataFrame\\[2pt]
    \rule{3cm}{0.3pt}\\
    get\_par\_id/nom()\\filtrer()\\get\_joueurs()
  };

  % Flèches composition (collection → élément)
  \draw[agg] (joueurs) -- node[right, font=\tiny] {contient} (joueur);
  \draw[agg] (matchs)  -- node[right, font=\tiny] {contient} (match);
  \draw[agg] (equipes) -- node[right, font=\tiny] {contient} (equipe);

  % Lien Equipe → Joueurs
  \draw[->, dashed, thick] (equipe.south) to[out=-90, in=90]
    node[right, font=\tiny\itshape, pos=0.5] {get\_joueurs()} (joueurs.north);
\end{tikzpicture}
\caption{Diagramme de classes simplifié du modèle objet.}
\label{fig:classes}
\end{figure}

\subsection{La classe Joueur}

La classe \classe{Joueur} encapsule une \code{pd.Series} représentant une ligne
du DataFrame joueurs. Ses propriétés principales sont :

\begin{itemize}
  \item \code{id}, \code{nom} : identifiant et nom du joueur.
  \item \code{naissance} : date de naissance sous forme de \code{pd.Timestamp}.
  \item \code{age} : calculé dynamiquement à partir de la date du jour.
  \item \code{equipe} : équipe d'appartenance (optionnel selon le sport).
  \item \code{taille} : taille en cm (optionnel selon le sport).
\end{itemize}

La classe \classe{Joueurs} est la collection correspondante. Elle supporte
l'itération Python (\code{\_\_iter\_\_}), l'indexation (\code{\_\_getitem\_\_})
et expose des méthodes statistiques : \code{age\_moyen()},
\code{taille\_moyenne()}, \code{resume()}.

\subsection{La classe Match}

\classe{Match} encapsule une ligne de DataFrame matchs. La propriété
\code{gagnant} applique la logique suivante (par ordre de priorité) :
\begin{enumerate}
  \item Si une colonne \code{col\_gagnant} est configurée, elle est utilisée
        directement.
  \item Sinon, le gagnant est déduit automatiquement en comparant
        \code{score\_1} et \code{score\_2}.
  \item En cas d'égalité ou de données manquantes, \code{None} est retourné.
\end{enumerate}

\classe{Matchs} expose notamment \code{bilan(equipe\_id)},
\code{classement()} et \code{score\_moyen()}, permettant des analyses
statistiques directement sur la collection.

\subsection{BaseLoader et pattern Template Method}

\classe{BaseLoader} est une classe abstraite (ABC Python) qui définit
l'infrastructure commune à tous les loaders. Elle implémente le
\textbf{patron de conception Template Method} : les sous-classes n'ont
qu'à déclarer \code{SPORT\_FOLDER} et à implémenter les méthodes
\code{load\_X()}, tout le reste (construction du chemin, lecture CSV, parsing
des dates) étant géré par la classe mère.

\begin{lstlisting}[caption={BaseLoader.py -- classe de base pour tous les loaders}]
class BaseLoader(ABC):
    _DATA_ROOT: Path = Path(__file__).parent.parent.parent / "data"
    SPORT_FOLDER: str = ""

    @property
    def root(self) -> Path:
        return self._DATA_ROOT / self.SPORT_FOLDER

    def _load_csv(self, filename, dtype=None, date_cols=None):
        df = pd.read_csv(self.root / filename, dtype=dtype or {})
        for col in (date_cols or []):
            if col in df.columns:
                df[col] = pd.to_datetime(df[col], errors="coerce")
        return df
\end{lstlisting}

% ------------------------------------------------------------
\section{Fonctionnalités de l'application}
\label{sec:fonct}
% ------------------------------------------------------------

\subsection{Interface utilisateur}

L'application s'exécute en ligne de commande via \code{python \_\_main\_\_.py}.
L'utilisateur navigue dans des menus hiérarchiques numérotés. La figure
ci-dessous schématise le flux de navigation principal.

\begin{figure}[h!]
\centering
\begin{tikzpicture}[
  box/.style={draw, rounded corners, fill=blue!8, font=\small,
              minimum width=3.5cm, minimum height=0.9cm, align=center},
  arr/.style={->, thick, >=Stealth}
]
  \node[box] (start) at (0, 0)    {Lancement\\(\_\_main\_\_.py)};
  \node[box] (menu1) at (0, -1.8) {Menu principal\\(choix du sport)};
  \node[box] (menu2) at (0, -3.6) {Menu du sport\\(équipe / joueurs / stats)};
  \node[box] (input) at (0, -5.4) {Saisie ou sélection\\paginée};
  \node[box] (result)at (0, -7.2) {Affichage du résultat};
  \node[box] (export)at (4, -7.2) {Export CSV};
  \node[box] (retour)at (0, -9.0) {Retour au menu};

  \draw[arr] (start)  -- (menu1);
  \draw[arr] (menu1)  -- (menu2);
  \draw[arr] (menu2)  -- (input);
  \draw[arr] (input)  -- (result);
  \draw[arr] (result) -- node[right, font=\footnotesize] {[e]} (export);
  \draw[arr] (result) -- (retour);
  \draw[arr] (retour.west) to[out=180, in=180, distance=2cm] (menu1.west);
\end{tikzpicture}
\caption{Flux de navigation de l'interface CLI.}
\label{fig:flux}
\end{figure}

\subsection{Sélecteur paginé}

Pour les fonctionnalités nécessitant de choisir parmi une longue liste (joueurs,
équipes, tournois), un \textbf{sélecteur paginé} a été implémenté. Il affiche
20 éléments par page et propose les commandes suivantes :

\begin{center}
\begin{tabular}{cl}
  \toprule
  \textbf{Touche} & \textbf{Action} \\
  \midrule
  \code{<} & Page précédente \\
  \code{>} & Page suivante \\
  \code{f} & Appliquer un filtre (catégorie, année, ou recherche textuelle) \\
  \code{r} & Réinitialiser les filtres \\
  numéro   & Sélectionner l'élément correspondant \\
  \code{0} & Annuler et revenir au menu \\
  \bottomrule
\end{tabular}
\end{center}

Le filtre de recherche supporte la \textbf{correspondance par préfixe de mot} :
une recherche sur \code{"leb"} renverra par exemple \textit{LeBron James}.

\subsection{Export CSV}

Après tout affichage de résultat tabulaire, l'utilisateur peut appuyer sur
\code{[e]} pour exporter le tableau au format CSV. Le fichier est créé dans
le répertoire courant avec un nom automatique :

\begin{center}
  \code{\{label\_nettoyé\}\_YYYYMMDD\_HHMMSS.csv}
\end{center}

L'encodage utilisé est \code{utf-8-sig} (compatible avec Microsoft Excel).

\subsection{Fonctionnalités par sport}

Le tableau~\ref{tab:fonct} liste les fonctionnalités disponibles pour chacun
des neuf sports.

\begin{center}
\begin{longtable}{lp{9cm}}
  \caption{Fonctionnalités par sport.}
  \label{tab:fonct}\\
  \toprule
  \textbf{Sport} & \textbf{Fonctionnalités} \\
  \midrule
  \endfirsthead
  \toprule
  \textbf{Sport} & \textbf{Fonctionnalités} \\
  \midrule
  \endhead
  \bottomrule
  \endfoot
  Basketball &
    Top équipes offensives · Stats d'une équipe · Roster ·
    Classement défensif · Fiche joueur · Liste joueurs · Classement par points \\
  Football UCL &
    Classement par points (groupes) · Bracket des phases finales ·
    Stats d'une équipe · Fiche joueur · Liste joueurs \\
  Football ligues &
    Classement par ligue · Stats d'une équipe · Fiche joueur · Liste joueurs \\
  League of Legends &
    Stats d'une équipe · Champions picks/bans · Roster ·
    Fiche joueur · Liste joueurs · Classement points \\
  CS2 &
    Stats d'une équipe · Classement · Roster ·
    Fiche joueur · Liste joueurs \\
  Tennis ATP/WTA &
    Classement par victoires · Stats d'un joueur ·
    Bracket d'un tournoi · Fiche joueur · Liste joueurs \\
  Volleyball &
    Classement (H/F) · Stats d'une équipe · Roster ·
    Fiche joueur · Liste joueurs \\
  StarCraft~II &
    Classement joueurs · Stats d'un joueur ·
    Bracket d'un tournoi · Fiche joueur · Liste joueurs \\
  Badminton &
    Classement · Stats d'un joueur · Bracket ·
    Fiche joueur · Liste joueurs \\
  Échecs &
    Classement · Stats d'un joueur · Bracket ·
    Fiche joueur · Liste joueurs \\
\end{longtable}
\end{center}

\subsection{Agenda multi-sport}

La fonction \code{agenda\_recents(sources, n)} est accessible depuis le menu
principal. Elle agrège les derniers matchs de tous les sports sélectionnés et
les affiche dans un tableau unifié :

\begin{center}
\begin{tabular}{cllcc}
  \toprule
  \textbf{Sport} & \textbf{Date} & \textbf{Équipe 1} & \textbf{Équipe 2} & \textbf{Score} \\
  \midrule
  Basketball & 2023-04-07 & Los Angeles Lakers & Boston Celtics & 115 -- 112 \\
  Tennis ATP  & 2024-01-28 & Jannik Sinner & Alexander Zverev & 2 -- 0 \\
  LoL        & 2024-03-15 & G2 Esports & Fnatic & 1 -- 0 \\
  \bottomrule
\end{tabular}
\end{center}

\subsection{Fonction générique \code{afficher\_classement}}

Toutes les fonctionnalités de classement reposent sur la même fonction
partagée \code{afficher\_classement()} du module \fichier{générique.py}.
Elle calcule les points selon la règle standard (3~pts victoire, 1~pt nul,
0~pt défaite) et supporte optionnellement l'affichage d'un séparateur
de qualification (\code{top\_qualifies}).

\subsection{Bracket de tournoi}

La fonction \code{afficher\_bracket()} affiche les confrontations tour par tour
sous forme ASCII, avec les gagnants mis en évidence en vert (code ANSI).
Elle supporte les compétitions à aller-retour (deux manches) via le paramètre
\code{deux\_manches}.

% ------------------------------------------------------------
\section{Tests et validation}
\label{sec:tests}
% ------------------------------------------------------------

\subsection{Organisation des tests}

Les tests sont organisés dans le répertoire \fichier{test/} selon la même
structure que le code source :

\begin{itemize}
  \item \fichier{test/Analysis/} --- 10 fichiers, un par module sport et
        pour \fichier{générique.py}.
  \item \fichier{test/Model/} --- 5 fichiers (Joueur, Joueurs, Match, Matchs,
        Equipe/Equipes).
  \item \fichier{test/Parsers/} --- 1 fichier (BaseLoader).
  \item \fichier{test/Common/} --- 1 fichier (utils.py).
\end{itemize}

\subsection{Exécution des tests}

Les tests s'exécutent via la commande :

\begin{lstlisting}[language=bash, numbers=none]
python -m pytest --cov=src --cov-report=term-missing
\end{lstlisting}

\subsection{Exemples de tests}

Les tests du modèle vérifient la logique métier indépendamment des données
réelles, en construisant des DataFrames de test minimaux :

\begin{lstlisting}[caption={Extrait de test pour la classe Matchs}]
import pandas as pd
from src.Model.match import Matchs

def test_bilan_equipe():
    df = pd.DataFrame({
        "equipe1": ["A", "A", "B"],
        "equipe2": ["B", "C", "C"],
        "score1":  [2,   1,   0],
        "score2":  [1,   2,   3],
    })
    matchs = Matchs(df, col_equipe1="equipe1", col_equipe2="equipe2",
                    col_score1="score1", col_score2="score2")
    bilan = matchs.bilan("A")
    assert bilan["victoires"] == 1
    assert bilan["defaites"] == 1
    assert bilan["joues"] == 2
\end{lstlisting}

\subsection{Formatage et linting}

L'ensemble du code est formaté avec \textbf{black} (version 25.1.0) et vérifié
avec \textbf{Flake8}. Le respect de ces outils garantit une cohérence stylistique
sur l'ensemble du projet.

% ------------------------------------------------------------
\section{Résultats et démonstration}
\label{sec:resultats}
% ------------------------------------------------------------

Cette section présente quelques captures d'écran illustrant le fonctionnement
de l'application.

\subsection{Menu principal}

\screenshot{Menu principal de l'application}{
  L'utilisateur sélectionne le sport désiré parmi les options numérotées.
  Une option « Agenda » permet de consulter les derniers matchs toutes
  disciplines confondues.
}

\subsection{Sélecteur de joueurs paginé}

\screenshot{Sélecteur paginé avec filtre de recherche}{
  Exemple de sélection d'un joueur NBA : la liste est paginée (20 éléments
  par page) et un filtre par préfixe permet de retrouver rapidement un joueur.
  La barre de navigation [<] Préc. | [>] Suiv. | [f] Filtrer | [r] Reset
  est visible en bas de l'écran.
}

\subsection{Fiche d'un joueur}

\screenshot{Fiche individuelle d'un joueur NBA}{
  La fiche d'un joueur NBA affiche ses informations complètes : nom, poste,
  numéro de maillot, date de naissance, nationalité et taille.
  L'option [e] Exporter en CSV est disponible après chaque résultat.
}

\subsection{Classement NBA}

\screenshot{Classement NBA par points (Regular Season 2022--2023)}{
  Le classement est calculé depuis les scores des matchs (3 pts victoire,
  1 pt nul, 0 pt défaite) et affiché sous forme de tableau trié.
}

\subsection{Bracket de tennis}

\screenshot{Bracket du tournoi de Roland-Garros 2024 (ATP)}{
  Les rounds sont affichés de gauche à droite (R128 → F).
  Le vainqueur de chaque match est mis en évidence en vert dans le terminal.
}

\subsection{Statistiques d'équipe LoL}

\screenshot{Statistiques de l'équipe G2 Esports (EMEA LEC)}{
  Les statistiques moyennes par match incluent : kills, assists, morts, or
  total, tourelles, dragons, barons et durée moyenne de la partie en minutes.
}

\subsection{Agenda multi-sport}

\screenshot{Agenda des 20 derniers matchs toutes disciplines confondues}{
  L'agenda agrège et trie chronologiquement les matchs de tous les sports
  disponibles. Chaque ligne indique le sport, la date, les deux équipes/joueurs
  et le score.
}

% ------------------------------------------------------------
\section{Conclusion}
\label{sec:conclusion}
% ------------------------------------------------------------

Ce projet a permis de développer une application complète de traitement de
données sportives couvrant dix disciplines, depuis l'ingestion de fichiers CSV
hétérogènes jusqu'à l'affichage interactif en ligne de commande.

\subsection{Bilan technique}

Du point de vue technique, le projet a atteint ses objectifs : la séparation
en couches indépendantes (Model, Parsers, Analysis, Common) garantit la
maintenabilité du code et facilite l'ajout de nouveaux sports. Le pattern
\textit{lazy loading} des loaders assure des performances satisfaisantes,
même avec des jeux de données volumineux. La couverture par tests unitaires
protège les composants critiques du modèle objet.

L'interface paginée avec filtres représente une amélioration significative de
l'expérience utilisateur, notamment pour les sports avec de nombreux joueurs
(NBA : 431 joueurs ; Tennis : plusieurs milliers de joueurs ATP/WTA).

\subsection{Difficultés rencontrées}

La principale difficulté a résidé dans l'\textbf{hétérogénéité des données} :
chaque source CSV utilise des conventions différentes (noms de colonnes, formats
de dates, représentation des scores). La mise en place d'une abstraction
suffisamment générique pour couvrir tous les sports --- tout en restant simple
à utiliser dans les modules d'analyse --- a nécessité plusieurs itérations
de conception.

La gestion des données manquantes (valeurs NaN dans certains jeux de données)
a également demandé une attention particulière, en particulier pour les calculs
statistiques et l'affichage.

\subsection{Perspectives d'amélioration}

Plusieurs extensions seraient envisageables :
\begin{itemize}
  \item Ajout d'une interface graphique (Tkinter ou interface web) pour
        améliorer l'accessibilité.
  \item Intégration de données en temps réel via des API sportives publiques.
  \item Implémentation de visualisations graphiques (histogrammes, courbes
        de progression) avec matplotlib.
  \item Extension à de nouveaux sports ou compétitions (Formule~1, rugby, etc.).
\end{itemize}

\bigskip

Ce projet de première année à l'ENSAI nous a permis de consolider nos
compétences en programmation Python orientée objet, en manipulation de données
avec pandas, et en génie logiciel (tests, versionnage, documentation). Il
illustre concrètement comment une architecture en couches permet de gérer la
complexité d'une application de taille réelle.

% ------------------------------------------------------------
% RÉFÉRENCES
% ------------------------------------------------------------
\newpage
\begin{thebibliography}{9}
  \bibitem{pandas}
    McKinney, W. (2010).
    \textit{Data Structures for Statistical Computing in Python}.
    Proceedings of the 9th Python in Science Conference.

  \bibitem{pep8}
    van Rossum, G., Warsaw, B., Coghlan, A. (2001).
    \textit{PEP 8 -- Style Guide for Python Code}.
    Python Software Foundation.

  \bibitem{pytest}
    Krekel, H. et al. (2004--2024).
    \textit{pytest: helps you write better programs}.
    \url{https://docs.pytest.org/}

  \bibitem{kaggle_nba}
    Jeu de données NBA 2022--2023.
    Kaggle (source publique).

  \bibitem{kaggle_tennis}
    Jeu de données Tennis ATP/WTA 2024.
    Jeff Sackmann, Tennis Abstract.
\end{thebibliography}

% ============================================================
% ANNEXES
% ============================================================
\newpage
\appendix
\section*{Annexes}
\addcontentsline{toc}{section}{Annexes}

% -----------------------------------------------
\subsection*{Annexe A --- Bilan personnel de Jade Colin}
\addcontentsline{toc}{subsection}{Annexe A --- Bilan personnel de Jade Colin}

\begin{tcolorbox}[colback=gray!5, colframe=gray!50, title=Bilan personnel]
  \textit{[À compléter par Jade Colin]}

  \medskip
  \textbf{Compétences acquises :} \\
  \ldots

  \medskip
  \textbf{Difficultés rencontrées :} \\
  \ldots

  \medskip
  \textbf{Ce que j'aurais fait différemment :} \\
  \ldots
\end{tcolorbox}

% -----------------------------------------------
\subsection*{Annexe B --- Bilan personnel de Mathias Douillard}
\addcontentsline{toc}{subsection}{Annexe B --- Bilan personnel de Mathias Douillard}

\begin{tcolorbox}[colback=gray!5, colframe=gray!50, title=Bilan personnel]
  \textit{[À compléter par Mathias Douillard]}

  \medskip
  \textbf{Compétences acquises :} \\
  \ldots

  \medskip
  \textbf{Difficultés rencontrées :} \\
  \ldots

  \medskip
  \textbf{Ce que j'aurais fait différemment :} \\
  \ldots
\end{tcolorbox}

% -----------------------------------------------
\subsection*{Annexe C --- Bilan personnel d'Adam [Nom]}
\addcontentsline{toc}{subsection}{Annexe C --- Bilan personnel d'Adam [Nom]}

\begin{tcolorbox}[colback=gray!5, colframe=gray!50, title=Bilan personnel]
  \textit{[À compléter par Adam [Nom]]}

  \medskip
  \textbf{Compétences acquises :} \\
  \ldots

  \medskip
  \textbf{Difficultés rencontrées :} \\
  \ldots

  \medskip
  \textbf{Ce que j'aurais fait différemment :} \\
  \ldots
\end{tcolorbox}

% -----------------------------------------------
\subsection*{Annexe D --- Bilan personnel de Mohamed Fouad Chaabane}
\addcontentsline{toc}{subsection}{Annexe D --- Bilan personnel de Mohamed Fouad Chaabane}

\begin{tcolorbox}[colback=gray!5, colframe=gray!50, title=Bilan personnel]
  \textit{[À compléter par Mohamed Fouad Chaabane]}

  \medskip
  \textbf{Compétences acquises :} \\
  \ldots

  \medskip
  \textbf{Difficultés rencontrées :} \\
  \ldots

  \medskip
  \textbf{Ce que j'aurais fait différemment :} \\
  \ldots
\end{tcolorbox}

% -----------------------------------------------
\subsection*{Annexe E --- Structure des répertoires}
\addcontentsline{toc}{subsection}{Annexe E --- Structure des répertoires}

\begin{Verbatim}[fontsize=\small, frame=single]
Projet-1A-2026/
├── __main__.py           # Point d'entrée CLI (~1400 lignes)
├── requirements.txt      # pandas, pytest, pytest-cov, black
├── README.md
├── AUTHORS.md
├── data/
│   ├── Basketball/       # game.csv, player.csv, team.csv
│   ├── tennis/           # atp_matches_2024.csv, wta_*
│   ├── Lol/              # match.csv, player.csv, coach.csv, team.csv
│   ├── counter_strike_2/ # match.csv, player.csv, coach.csv, team.csv
│   ├── volleyball/       # match_men/women.csv, player_men/women.csv
│   ├── football_champions_league/
│   ├── Football/
│   ├── starcraft_2/
│   ├── badminton/
│   └── chess/
├── src/
│   ├── Model/            # Joueur, Match, Equipe, Sport
│   ├── Parsers/          # BaseLoader + 10 loaders
│   ├── Analysis/         # générique.py + 9 modules sport
│   └── Common/           # utils.py
└── test/
    ├── Analysis/         # 10 fichiers de tests
    ├── Model/            # 5 fichiers de tests
    ├── Parsers/          # 1 fichier de tests
    └── Common/           # 1 fichier de tests
\end{Verbatim}

\end{document}
